\chapter{Strings in Python: Working with Text}

\section{Introduction: Real-Life Hook}

Students work with text every day:  

\begin{itemize}
    \item Writing names and addresses  
    \item Sending messages to friends  
    \item Storing favorite quotes or book titles  
\end{itemize}

In Python, \textbf{strings} let us handle all these text tasks easily.

\section{Concept: What is a String?}

\begin{definition}{String}
A \textbf{string} is a sequence of characters enclosed in quotes, used to store text in Python.
\end{definition}

\subsection{Creating Strings}

\begin{verbatim}
name = "Amina"
message = 'Hello, how are you?'
\end{verbatim}

- Strings can use single or double quotes  
- Can include letters, numbers, spaces, and symbols  

\subsection{Accessing Characters}

\begin{verbatim}
name = "Amina"
print(name[0])  # 'A'
print(name[2])  # 'i'
\end{verbatim}

- Indexing starts at 0  
- Negative indices count from the end: \texttt{name[-1]} is 'a'

\subsection{String Slicing}

\begin{verbatim}
text = "Python"
print(text[0:4])   # 'Pyth'
print(text[2:])    # 'thon'
print(text[:3])    # 'Pyt'
\end{verbatim}

- \texttt{[start:end]} extracts a substring from start index up to (but not including) end index  

\section{String Operations}

\begin{itemize}
    \item Concatenation: joining strings using \texttt{+}  

\begin{verbatim}
first = "Hello"
second = "World"
print(first + " " + second)  # 'Hello World'
\end{verbatim}

    \item Repetition: repeat strings using \texttt{*}  

\begin{verbatim}
print("Ha" * 3)  # 'HaHaHa'
\end{verbatim}

    \item Length: find string length using \texttt{len()}  

\begin{verbatim}
text = "Python"
print(len(text))  # 6
\end{verbatim}
\end{itemize}

\section{String Methods}

\begin{verbatim}
text = "python programming"
print(text.upper())    # 'PYTHON PROGRAMMING'
print(text.lower())    # 'python programming'
print(text.title())    # 'Python Programming'
print(text.replace("python", "Java"))  # 'Java programming'
\end{verbatim}

\subsection{Checking Text}

\begin{verbatim}
text = "12345"
print(text.isdigit())  # True
print(text.isalpha())  # False
\end{verbatim}

\section{Looping Through Strings}

\begin{verbatim}
word = "Python"
for char in word:
    print(char)
\end{verbatim}

Output:

\begin{verbatim}
P
y
t
h
o
n
\end{verbatim}

\section{Mini-Project: Word Analyzer}

Create a program that asks the user to enter a word and prints:  

\begin{itemize}
    \item Number of characters  
    \item Uppercase version  
    \item First and last character  
\end{itemize}

\begin{verbatim}
word = input("Enter a word: ")
print("Length:", len(word))
print("Uppercase:", word.upper())
print("First character:", word[0])
print("Last character:", word[-1])
\end{verbatim}

\section{Step-by-Step Diagram: String Operations Flow}

\begin{center}
\begin{tikzpicture}[node distance=1.5cm]
\node (start) [draw, ellipse, fill=green!20] {Start Program};
\node (input) [draw, rectangle, fill=yellow!20, below of=start] {Get String Input};
\node (process) [draw, rectangle, fill=orange!20, below of=input] {Apply String Operations / Methods};
\node (output) [draw, ellipse, fill=green!20, below of=process] {Display Results};
\draw[->] (start) -- (input);
\draw[->] (input) -- (process);
\draw[->] (process) -- (output);
\end{tikzpicture}
\end{center}

\section{Key Concepts}

\begin{itemize}
    \item Strings store text in quotes and are indexed starting from 0  
    \item Strings can be sliced, concatenated, repeated, and measured  
    \item Methods allow modification, replacement, and checking of strings  
    \item Loops can iterate through each character in a string  
\end{itemize}

\section{End-of-Chapter Exercises}

\begin{exercise}{Very Short Questions}
\begin{enumerate}
    \item What is a string in Python?  
    \item How do you find the length of a string?  
    \item Name two string methods to change case.  
\end{enumerate}
\end{exercise}

\begin{exercise}{Short Answer Questions}
\begin{enumerate}
    \item Write a program to print the first and last character of a word.  
    \item How do you check if a string contains only numbers?  
    \item Explain string slicing with an example.  
\end{enumerate}
\end{exercise}

\begin{exercise}{Activity-Based Questions}
\begin{enumerate}
    \item Ask the user to enter their full name and print it in uppercase and lowercase.  
    \item Count the number of vowels in a word entered by the user.  
    \item Create a program to replace all spaces in a sentence with hyphens (``-'').  
\end{enumerate}
\end{exercise}

\section{Teacher Notes}

\marginnote{
\textbf{Learning Objectives:}  
\begin{itemize}
    \item Students understand what strings are and how to manipulate them  
    \item Students can apply string methods and slicing  
    \item Students can loop through strings for character-level operations  
\end{itemize}

\textbf{Teaching Tips:}  
\begin{itemize}
    \item Use real-life examples like names, messages, or sentences  
    \item Demonstrate indexing, slicing, and methods live in class  
\end{itemize}

\textbf{Assessment:}  
\begin{itemize}
    \item Check correctness of string operations  
    \item Observe understanding of indexing, slicing, and loops over strings  
\end{itemize}
}